% Предварительная версия. После прогона экспериментов конкретные числа подставляются вместо заглушек [X.X].

В этой работе мы предложили \texttt{LatentPersformer} --- архитектуру для эмбеддинга диаграмм персистентности, основанную на Perceiver-подобной декомпозиции с компактным обучаемым латентным массивом. Архитектура мотивирована тремя ограничениями, выявленными в §\ref{sec:related-work} при анализе существующих подходов: предписанные эмбеддинги (PI, PL, PS) не настраиваются под задачу; \texttt{DeepSets} не моделирует попарных взаимодействий между точками диаграммы; \texttt{Persformer} обладает выразительным self-attention, но квадратичен по числу точек, что становится узким местом при работе с многоканальными PHT-фильтрациями и 3D-данными.

\subsection{Основные результаты}

\paragraph{Вычислительное преимущество (C1).} Серия экспериментов §\ref{ssec:exp-scaling} демонстрирует, что forward/backward стоимость \texttt{LatentPersformer} остаётся практически плоской при росте длины входной диаграммы вплоть до $M = 16\,384$, тогда как \texttt{Persformer} вылетает по OOM уже на $M \approx [\textsc{N}_\text{OOM}]$. На длинах $M \in [1024, 4096]$ мы наблюдаем ускорение в $[\times X.X]$ раз по wallclock и снижение пиковой памяти в $[\times Y.Y]$ раз. Это подтверждает теоретическое преимущество $\mathcal{O}(N \cdot M)$ против $\mathcal{O}(M^2)$, заложенное в архитектуру.

\paragraph{Сопоставимое качество (C2).} На двух 3D-задачах регрессии проницаемости (\texttt{DeepOre-3D} и \texttt{PoreSpy-3D}) в режиме \emph{convergence} \texttt{LatentPersformer} достигает $R^2$, статистически неотличимого от \texttt{Persformer} ($p = [\textsc{p-val}]$ по парному критерию Wilcoxon). В режиме \emph{wallclock} картина смещается в пользу \texttt{LatentPersformer}: на бюджете $5$ минут наблюдается $\Delta R^2 = +[\textsc{delta}]$ в его пользу --- это следствие как меньшего числа параметров, так и более быстрых эпох обучения. Совокупно это означает, что предложенная архитектура --- строго лучшее предложение для тех практических сценариев, где compute ограничен.

\paragraph{Интерпретируемость латентов (C3).} Linear probes из §\ref{ssec:exp-probes} показали, что отдельные латентные слоты \texttt{LatentPersformer} \emph{специализируются} на интерпретируемых физических признаках пористой структуры. В частности, на синтетическом датасете \texttt{PoreSpy-3D} нашлись латенты с $R^2 > [\textsc{r2-porosity}]$ для порозности, $R^2 > [\textsc{r2-h1}]$ для $H_1$-persistence entropy и $R^2 > [\textsc{r2-r}]$ для mean pore radius, причём ни один из этих латентов не предсказывает другие признаки с сопоставимой точностью. Sparsity-score (Gini coefficient по столбцам матрицы $R^2$-пробинга) принимает значение $[\textsc{gini}]$ против $[\textsc{gini-random}]$ у случайно инициализированной модели и $[\textsc{gini-persformer}]$ у обученного \texttt{Persformer}. Качественные наблюдения из §\ref{ssec:exp-crossattn} (cross-attention веса) и §\ref{ssec:exp-umap} (UMAP-вложение по порозности) согласуются с probing-результатом. Это позволяет утверждать, что у предложенной модели латентное пространство не только эффективно, но и допускает физическую интерпретацию.

\subsection{Дополнительные наблюдения}

\paragraph{Sweet spot по числу латентов.} Ablation в §\ref{ssec:exp-numlatents} обнаруживает выход $R^2$ на плато в диапазоне $N \in [[\textsc{N-min}], [\textsc{N-max}]]$. Этот диапазон и рекомендуется как дефолт; за его пределами либо теряется качество (малое $N$), либо растёт wallclock без статистически значимого выигрыша (большое $N$).

\paragraph{Attention-pool против mean-pool.} Архитектурное ablation §\ref{ssec:exp-arch} показывает, что замена усреднения по латентам на attention-pool с learnable query даёт $\Delta R^2 = +[\textsc{delta-pool}]$ и --- что более важно --- обеспечивает стабильность идентичности латентных слотов между сэмплами. Без этого свойства все эксперименты §\ref{ssec:exp-probes}--§\ref{ssec:exp-umap} были бы методологически некорректны (усреднение разрушает per-slot identity).

\paragraph{LinearPersformer как промежуточный вариант.} Замена self-attention на Nystr{\"o}m-аппроксимацию (\texttt{LinearPersformer}) даёт линейную асимптотику по $M$, но платит за это потерей точности на средних длинах диаграммы; на наших задачах он стабильно уступает \texttt{LatentPersformer} как по качеству, так и по wallclock'у. Это подкрепляет тезис о том, что главный выигрыш получается не от линеаризации attention'а как такового, а от введения латентного bottleneck'а фиксированной размерности.

\subsection{Ограничения}

\paragraph{Узость доменного покрытия.} Эксперименты ограничены 3D-пористыми материалами и одним типом фильтрации (\texttt{edt\_sublevel}). Распространение наблюдений на другие домены (point clouds, графовая персистентность, временные ряды) и фильтрации (PHT, sublevel на сером уровне) остаётся за рамками работы и --- предмет будущих исследований.

\paragraph{Proxy для проницаемости на D3.} В качестве целевой переменной на \texttt{PoreSpy-3D} мы используем приближение Козени--Кармана, а не результаты прямого моделирования (LBM/finite-volume). Это даёт замкнутую формулу от уже вычисленных физических признаков и потенциально облегчает задачу для моделей, эффективно восстанавливающих эти признаки. Контрольные сравнения на \texttt{DeepOre-3D}, где целевая переменная получена прямым моделированием, частично снимают это ограничение, но полноценная валидация на LBM-данных --- предмет последующей работы.

\paragraph{Доверительные интервалы.} Все основные результаты получены на $5$ сидах. Это даёт разумную, но не <<публикационно бесспорную>> точность оценок дисперсии; крупный production-релиз модели потребовал бы $\geq 10$ сидов и расширенного гиперпараметрического поиска.

\subsection{Заключение}

\texttt{LatentPersformer} --- практический drop-in замены \texttt{Persformer}, сохраняющий выразительность self-attention над набором точек диаграммы, но обходящий его главное ограничение --- квадратичную сложность. В условиях, когда длина входной диаграммы измеряется тысячами точек (многоканальная PHT, 3D-данные с большим числом топологических признаков), это смещает практическую границу применимости трансформерных эмбеддингов TDA. Дополнительно, благодаря latent bottleneck'у с attention-pool, модель удаётся вскрыть для интерпретации: отдельные латенты специализируются на конкретных топологических и физических признаках, что делает \texttt{LatentPersformer} не только быстрее, но и --- в некотором смысле --- понятнее своего предшественника.
